\section{To-Do}
\label{sec:todo}

The most pressing extensions to the model can be ranked by which would most directly test or change the qualitative result. Three stand out, and each is straightforward to implement on top of the existing simulator.

The first is standing variation with a fitness cost for R outside the edge. This is the extension most likely to flip the dormant-loses-everywhere result. The simulator currently starts every run with $\Nzero$ wild-type cells and no resistant cells, and treats R as neutral in the core. The minimal change is two parameters: a pre-treatment phase of duration $T_{\mathrm{pre}}$ during which the antibiotic is absent ($\bEdge = \dEdge = 1$ everywhere) and a fitness cost $c_R \in (0, 1)$ that reduces R cells' birth rate to $b \cdot (1 - c_R)$ throughout the pre-treatment phase, where $b$ is the compartment-appropriate birth rate. During $T_{\mathrm{pre}}$, mutations arise at rate $\mu$ per wild-type birth, R cells incur the fitness cost wherever they are alive, and a standing R frequency builds up that depends on $\mu$, $c_R$, and the spatial structure. At $t = T_{\mathrm{pre}}$ the antibiotic switches on, current dynamics resume, and rescue is measured as before. The expected output: in the well-mixed population, standing R is selected down to a low frequency of order $\mu / c_R$. In the dormant biofilm, R mutations produced in the edge can be pushed into the core via the conveyor belt, where there is no cell turnover and therefore no selection against them; over time, the dormant core preserves a reservoir of R cells that would have been culled in a well-mixed population. When treatment starts, the dormant biofilm enters with this preserved reservoir, and may rescue more often than the well-mixed reference at low dose, where our current model says it always loses. The natural sweep is $c_R$ across a small range with $T_{\mathrm{pre}}$ long enough for the standing R frequency to equilibrate.

The second is the antibiotic's mode of action. The simulator's $\dEdge$ is currently the only dose-sensitive parameter; $\bEdge$ stays at $1$. The minimal change is to make both edge rates dose-sensitive, with the choice of which actually varies selecting the drug's mode. In a purely bactericidal regime (current), the dose maps to $\dEdge$ alone and $\bEdge$ stays at $1$. In a purely bacteriostatic regime, the dose maps to $\bEdge$ instead, with $\bEdge$ shrinking from $1$ toward zero as dose increases, and $\dEdge$ remaining at $1$. The boundary speed is $\sE \cdot l$ with $\sE = \dEdge - \bEdge$ in either case. In the bactericidal case $\sE$ grows linearly with dose; in the bacteriostatic case $\sE$ saturates near $1$ as $\bEdge \to 0$, so the boundary speed has a hard ceiling no matter how high the dose is pushed. The race from Section~\ref{sec:reachingedge} would therefore have a fundamentally different high-dose limit: the boundary cannot run away from drift as quickly under a bacteriostatic drug as it can under a bactericidal one. The natural experiment is to run the same dose range under each mode and compare crossover doses, advantage magnitudes, and whether the dose-dependent flip survives in shape or is replaced by a saturating profile at high dose.

The third is dispersal as an in-simulation event. Currently the biofilm and well-mixed populations are simulated as separate runs. The minimal change is to allow a transition from biofilm to well-mixed at a specified time $T_{\mathrm{disp}}$ within a single run, by setting $l \to N(T_{\mathrm{disp}})$ at that moment so that every remaining cell is reclassified as edge. After $T_{\mathrm{disp}}$, the simulation continues with no core. The rescue probability of such a hybrid run is a direct in-model prediction of what a dispersal intervention does over the course of a treatment, including the dynamics of the transition itself. The natural sweep is $T_{\mathrm{disp}}$ across the dose response: intervening very early (close to $t = 0$) should produce dynamics close to a pure well-mixed run from the start; intervening very late (after most rescue events have occurred) should produce dynamics close to a pure biofilm run. The interesting regime is the intermediate one, and whether there is a $T_{\mathrm{disp}}$ that minimizes rescue probability at a given dose, which would correspond to the optimal timing of a dispersal intervention.
